\shead{Writing a number as a percentage -- and back}
\begin{knowbox}\begin{itemize}[leftmargin=1.4em,itemsep=3pt]
\item The sign $\%$ means ``out of $100$'': $\ 1\%=\dfrac1{100}=0.01$. \ So $\ a=a\times100\%$.
\item \textbf{Number $\to$ percentage:} multiply the number by $100$ and write the sign $\%$ after it.\par
Examples: $\ 3=3\times100\%=300\%;\quad 0.425=42.5\%;\quad \dfrac34=\dfrac34\times100\%=75\%$.
\item \textbf{Percentage $\to$ decimal:} divide by $100$ (move the decimal point two places to the left).\par
Examples: $\ 36\%=36:100=0.36;\quad 4.5\%=0.045;\quad 120\%=1.2$.
\item \textbf{Percentage $\to$ fraction:} write it as a fraction with denominator $100$, then simplify.\par
Example: $\ 85\%=\dfrac{85}{100}=\dfrac{17}{20}$.
\end{itemize}
\note{The sign $\%$ simply takes the place of ``divided by $100$''.}
\end{knowbox}
\gap
\shead{Calculating with percentages}
\begin{knowbox}\begin{itemize}[leftmargin=1.4em,itemsep=3pt]
\item \textbf{Adding and subtracting:} add or subtract the numbers and keep the sign $\%$:\par
$a\%+b\%=(a+b)\%;\qquad a\%-b\%=(a-b)\%$.
\item \textbf{Multiplying or dividing by an ordinary number $k$:} $\ a\%\times k=(a\times k)\%;\quad a\%:k=(a:k)\%$.
\item \textbf{Percentage times percentage:} $\ a\%\times b\%=\dfrac a{100}\times\dfrac b{100}$. \ Change both to decimals first:\par
$30\%\times40\%=0.3\times0.4=0.12=12\%$ \ (\emph{not} $1\,200\%$!).
\item \textbf{Percentage divided by percentage:} the two signs $\%$ cancel and the answer is an ordinary number: $\ a\%:b\%=a:b$. \ Example: $\ 45\%:5\%=9$.
\end{itemize}
\note{When in doubt, treat each percentage as a fraction with denominator $100$ (or as a decimal) and calculate as usual.}
\end{knowbox}
\gap
\shead{Expressions with several kinds of numbers}
\begin{knowbox}\begin{itemize}[leftmargin=1.4em,itemsep=3pt]
\item An expression may contain decimals, fractions, mixed numbers and percentages together. First change everything into \textbf{one kind} of number (all decimals, or all fractions), then follow the usual order of operations: brackets first, then $\times$ and $:$, then $+$ and $-$.
\item To find $x$, work backwards exactly as with whole numbers, after changing each percentage into a decimal or a fraction.
\end{itemize}
\end{knowbox}
\gap
\shead{Percentages in word problems}
\begin{knowbox}\begin{itemize}[leftmargin=1.4em,itemsep=3pt]
\item ``$A$ is $a\%$ of $B$'' means $A=\dfrac a{100}$ of $B$. \ Simplify the fraction to see the parts:\par
``$A$ is $40\%$ of $B$'' $\ \Rightarrow\ A=\dfrac25$ of $B$: \ $A$ is $2$ parts, $B$ is $5$ equal parts.
\item Then use the familiar methods with a bar model: \textbf{sum and ratio}, \textbf{difference and ratio}.
\item The whole of something is $100\%$. \ If parts of it are given as percentages, the part that is left is $100\%$ minus the given percentages.
\end{itemize}
\end{knowbox}
\gap
\begin{exbox}{\ExWord\ 1}
\pa{a} Write as percentages: $\ 0.36;\quad \dfrac7{20};\quad 3;\quad 1.075$.\par
\pa{b} Write as a decimal and as a fraction in simplest form: $\ 45\%;\quad 6.5\%;\quad 160\%$.
\guide
\pa{a} $0.36=0.36\times100\%=36\%$.\par
$\dfrac7{20}=\dfrac{7\times5}{20\times5}=\dfrac{35}{100}=35\%$ \ \ or \ \ $\dfrac7{20}=7:20=0.35=35\%$.\par
$3=3\times100\%=300\%$; \qquad $1.075=107.5\%$.\par
\pa{b} $45\%=45:100=0.45$; \quad $45\%=\dfrac{45}{100}=\dfrac9{20}$.\par
$6.5\%=6.5:100=0.065$; \quad $6.5\%=\dfrac{65}{1\,000}=\dfrac{13}{200}$.\par
$160\%=1.6$; \quad $160\%=\dfrac{160}{100}=\dfrac85$.
\end{exbox}
\gap
\begin{exbox}{\ExWord\ 2}
Calculate.
\begin{cols3}
\pa{a} $37\%+48\%$ & \pa{b} $126\%-58\%$ & \pa{c} $15.5\%\times6$\\
\pa{d} $40\%\times15\%$ & \pa{e} $72\%:9\%$ & \pa{g} $72\%:9$
\end{cols3}
\guide
\pa{a} \emph{Method 1.} $37\%+48\%=(37+48)\%=85\%$.\par
\emph{Method 2.} $37\%+48\%=0.37+0.48=0.85=85\%$.\par
\pa{b} $126\%-58\%=(126-58)\%=68\%$.\par
\pa{c} $15.5\%\times6=(15.5\times6)\%=93\%$.\par
\pa{d} $40\%\times15\%=0.4\times0.15=0.06=6\%$.\par
\pa{e} $72\%:9\%=0.72:0.09=8$.\par
\pa{g} $72\%:9=(72:9)\%=8\%$.
\note{Be careful: $40\%\times15\%$ is \emph{not} $600\%$, and $72\%:9\%$ is the number $8$, not $8\%$.}
\end{exbox}
\gap
\begin{exbox}{\ExWord\ 3}
Calculate: $\ A=4.5-1\dfrac34+63\%:18\%$.
\guide
\emph{Method 1 (change to decimals).}\par
$A=4.5-1.75+0.63:0.18=4.5-1.75+3.5=2.75+3.5=6.25$.\par
\emph{Method 2 (change to fractions).}\par
$A=\dfrac92-\dfrac74+\dfrac{63}{100}:\dfrac{18}{100}=\dfrac92-\dfrac74+\dfrac72=\dfrac{18}4-\dfrac74+\dfrac{14}4=\dfrac{25}4=6\dfrac14$.
\note{Both answers $A=6.25$ and $A=6\dfrac14$ are correct.}
\end{exbox}
\gap
\begin{exbox}{\ExWord\ 4}
Find $x$: $\ x:25\%+1.6=9:0.5$.
\guide
Work out the right-hand side: $9:0.5=18$.\par
So $x:25\%+1.6=18$, \ hence $x:25\%=18-1.6=16.4$.\par
Then $x=16.4\times25\%=16.4\times0.25=4.1$.\par
\emph{Check:} $4.1:0.25+1.6=16.4+1.6=18$.
\dapso{$x=4.1$.}
\end{exbox}
\gap
\begin{exbox}{\ExWord\ 5}
Lan and Hoa have $72$ stickers altogether. Hoa's stickers are $80\%$ of Lan's. How many stickers does each girl have?
\guide
$80\%=\dfrac{80}{100}=\dfrac45$, so Hoa's stickers are $\dfrac45$ of Lan's.
\begin{barfigC}
\rowlab{-0.15}{0.6}{Lan:}
\foreach \i in {0,...,4}{\blk{\i*0.8}{\i*0.8+0.8}{0.6}{0.42}{tmBlue}{}}
\rowlab{-0.15}{0}{Hoa:}
\foreach \i in {0,...,3}{\blk{\i*0.8}{\i*0.8+0.8}{0}{0.42}{tmGreen}{}}
\vbrace{4.1}{0.81}{-0.21}{$72$ stickers}
\end{barfigC}
Total number of equal parts: $5+4=9$.\par
One part: $72:9=8$ (stickers).\par
Lan: $8\times5=40$ stickers; \ Hoa: $72-40=32$ stickers.
\dapso{Lan $40$ stickers, Hoa $32$ stickers.}
\end{exbox}
\dividerbar{Practice}
\begin{probbox}{1}{\relax}
Write as percentages.
\begin{cols3}
\pa{a} $0.47$ & \pa{b} $1.36$ & \pa{c} $0.209$\\
\pa{d} $0.8$ & \pa{e} $2.5$ & \pa{g} $0.035$
\end{cols3}
\end{probbox}
\probsol{
Multiply by $100$ and write the sign $\%$:\par
\pa{a} $0.47=47\%$; \quad \pa{b} $1.36=136\%$; \quad \pa{c} $0.209=20.9\%$;\par
\pa{d} $0.8=80\%$; \quad \pa{e} $2.5=250\%$; \quad \pa{g} $0.035=3.5\%$.
\dapso{$47\%$; $136\%$; $20.9\%$; $80\%$; $250\%$; $3.5\%$.}}
\gap
\begin{probbox}{2}{\relax}
Write as percentages.
\begin{cols3}
\pa{a} $\dfrac35$ & \pa{b} $\dfrac9{20}$ & \pa{c} $\dfrac{21}{30}$\\
\pa{d} $\dfrac78$ & \pa{e} $\dfrac{13}4$ & \pa{g} $\dfrac{33}{40}$
\end{cols3}
\end{probbox}
\probsol{
\pa{a} $\dfrac35=\dfrac{60}{100}=60\%$. \quad \pa{b} $\dfrac9{20}=\dfrac{45}{100}=45\%$. \quad \pa{c} $\dfrac{21}{30}=\dfrac7{10}=\dfrac{70}{100}=70\%$.\par
\pa{d} $\dfrac78=7:8=0.875=87.5\%$. \quad \pa{e} $\dfrac{13}4=\dfrac{325}{100}=325\%$.\par
\pa{g} $\dfrac{33}{40}=33:40=0.825=82.5\%$.
\dapso{$60\%$; $45\%$; $70\%$; $87.5\%$; $325\%$; $82.5\%$.}}
\gap
\begin{probbox}{3}{\relax}
Write as decimals.
\begin{cols3}
\pa{a} $38\%$ & \pa{b} $4\%$ & \pa{c} $215\%$\\
\pa{d} $9.6\%$ & \pa{e} $300\%$ & \pa{g} $0.5\%$
\end{cols3}
\end{probbox}
\probsol{
Divide by $100$:\par
\pa{a} $38\%=0.38$; \quad \pa{b} $4\%=0.04$; \quad \pa{c} $215\%=2.15$;\par
\pa{d} $9.6\%=0.096$; \quad \pa{e} $300\%=3$; \quad \pa{g} $0.5\%=0.005$.
\dapso{$0.38$; $0.04$; $2.15$; $0.096$; $3$; $0.005$.}}
\gap
\begin{probbox}{4}{\relax}
Write as fractions in simplest form.
\begin{cols3}
\pa{a} $40\%$ & \pa{b} $85\%$ & \pa{c} $64\%$\\
\pa{d} $150\%$ & \pa{e} $2.5\%$ & \pa{g} $62.5\%$
\end{cols3}
\end{probbox}
\probsol{
\pa{a} $40\%=\dfrac{40}{100}=\dfrac25$. \quad \pa{b} $85\%=\dfrac{85}{100}=\dfrac{17}{20}$. \quad \pa{c} $64\%=\dfrac{64}{100}=\dfrac{16}{25}$.\par
\pa{d} $150\%=\dfrac{150}{100}=\dfrac32$. \quad \pa{e} $2.5\%=\dfrac{25}{1\,000}=\dfrac1{40}$. \quad \pa{g} $62.5\%=\dfrac{625}{1\,000}=\dfrac58$.
\dapso{$\dfrac25$; $\dfrac{17}{20}$; $\dfrac{16}{25}$; $\dfrac32$; $\dfrac1{40}$; $\dfrac58$.}}
\gap
\begin{probbox}{5}{\relax}
Compare (write $>$, $<$ or $=$).
\begin{cols2}
\pa{a} $0.7\ \ \dots\ \ 69\%$ & \pa{b} $\dfrac38\ \ \dots\ \ 37.5\%$\\
\pa{c} $125\%\ \ \dots\ \ 1\dfrac15$ & \pa{d} $\dfrac23\ \ \dots\ \ 66\%$
\end{cols2}
\end{probbox}
\probsol{
\pa{a} $0.7=70\%$ and $70\%>69\%$, so $0.7>69\%$.\par
\pa{b} $\dfrac38=3:8=0.375=37.5\%$, so $\dfrac38=37.5\%$.\par
\pa{c} $125\%=1.25$ and $1\dfrac15=1.2$; \ $1.25>1.2$, so $125\%>1\dfrac15$.\par
\pa{d} $66\%=\dfrac{66}{100}=\dfrac{33}{50}$, \ and $\dfrac23=\dfrac{100}{150}>\dfrac{99}{150}=\dfrac{33}{50}$, so $\dfrac23>66\%$.
\dapso{\pa{a} $>$;\ \ \pa{b} $=$;\ \ \pa{c} $>$;\ \ \pa{d} $>$.}}
\gap
\begin{probbox}{6}{\relax}
Write in order from smallest to largest: $\ 0.43,\ \ \dfrac25,\ \ 41.5\%,\ \ \dfrac9{20},\ \ 39\%$.
\end{probbox}
\probsol{
Change all of them into decimals:\par
$0.43$; \quad $\dfrac25=0.4$; \quad $41.5\%=0.415$; \quad $\dfrac9{20}=0.45$; \quad $39\%=0.39$.\par
Order: $0.39<0.4<0.415<0.43<0.45$.
\dapso{$39\%,\ \dfrac25,\ 41.5\%,\ 0.43,\ \dfrac9{20}$.}}
\gap
\begin{probbox}{7}{\relax}
Calculate.
\begin{cols2}
\pa{a} $28\%+56\%$ & \pa{b} $143\%-67\%$\\
\pa{c} $12.5\%+7.25\%+30.25\%$ & \pa{d} $90\%-45\%+15\%$
\end{cols2}
\end{probbox}
\probsol{
\pa{a} $28\%+56\%=(28+56)\%=84\%$.\par
\pa{b} $143\%-67\%=(143-67)\%=76\%$.\par
\pa{c} $12.5\%+7.25\%+30.25\%=(12.5+7.25+30.25)\%=50\%$.\par
\pa{d} $90\%-45\%+15\%=(90-45+15)\%=60\%$.
\dapso{$84\%$; $76\%$; $50\%$; $60\%$.}}
\gap
\begin{probbox}{8}{\relax}
Calculate.
\begin{cols2}
\pa{a} $23.4\%\times5$ & \pa{b} $252\%:12$\\
\pa{c} $6.8\%\times15$ & \pa{d} $99\%:4.5$
\end{cols2}
\end{probbox}
\probsol{
\pa{a} $23.4\%\times5=(23.4\times5)\%=117\%$.\par
\pa{b} $252\%:12=(252:12)\%=21\%$.\par
\pa{c} $6.8\%\times15=(6.8\times15)\%=102\%$.\par
\pa{d} $99\%:4.5=(99:4.5)\%=22\%$.
\dapso{$117\%$; $21\%$; $102\%$; $22\%$.}}
\gap
\begin{probbox}{9}{\relax}
Calculate and write the result as a decimal.
\begin{cols3}
\pa{a} $60\%\times35\%$ & \pa{b} $140\%\times25\%$ & \pa{c} $3.5\%\times80\%$\\
\pa{d} $84\%:12\%$ & \pa{e} $45\%:0.5\%$ & \pa{g} $7.2\%:16\%$
\end{cols3}
\end{probbox}
\probsol{
\pa{a} $60\%\times35\%=0.6\times0.35=0.21$.\par
\pa{b} $140\%\times25\%=1.4\times0.25=0.35$.\par
\pa{c} $3.5\%\times80\%=0.035\times0.8=0.028$.\par
\pa{d} $84\%:12\%=84:12=7$.\par
\pa{e} $45\%:0.5\%=45:0.5=90$.\par
\pa{g} $7.2\%:16\%=7.2:16=0.45$.
\dapso{$0.21$; $0.35$; $0.028$; $7$; $90$; $0.45$.}}
\gap
\begin{probbox}{10}{\relax}
Calculate.
\begin{cols2}
\pa{a} $2.4+35\%-\dfrac34$ & \pa{b} $1\dfrac12\times40\%+2.8:70\%$
\end{cols2}
\end{probbox}
\probsol{
\pa{a} $2.4+35\%-\dfrac34=2.4+0.35-0.75=2.75-0.75=2$.\par
\pa{b} $1\dfrac12\times40\%+2.8:70\%=1.5\times0.4+2.8:0.7=0.6+4=4.6$.
\dapso{\pa{a} $2$;\ \ \pa{b} $4.6$.}}
\gap
\begin{probbox}{11}{\relax}
Calculate.
\begin{cols2}
\pa{a} $25\%\times8+13.5:4.5-2\dfrac25$ & \pa{b} $4.8-1\dfrac14+36\%:144\%$
\end{cols2}
\end{probbox}
\probsol{
\pa{a} $25\%\times8+13.5:4.5-2\dfrac25=0.25\times8+3-2.4=2+3-2.4=2.6$.\par
\pa{b} $4.8-1\dfrac14+36\%:144\%=4.8-1.25+36:144=4.8-1.25+0.25=3.8$.
\dapso{\pa{a} $2.6$;\ \ \pa{b} $3.8$.}}
\gap
\begin{probbox}{12}{\relax}
Calculate in a clever way.
\begin{cols3}
\pa{a} $37\%\times2.5+63\%\times2.5$ & \pa{b} $125\%\times4.7\times80\%$ & \pa{c} $6.4\times45\%+6.4\times55\%$
\end{cols3}
\end{probbox}
\probsol{
\pa{a} $37\%\times2.5+63\%\times2.5=(37\%+63\%)\times2.5=100\%\times2.5=1\times2.5=2.5$.\par
\pa{b} $125\%\times4.7\times80\%=(1.25\times0.8)\times4.7=1\times4.7=4.7$.\par
\pa{c} $6.4\times45\%+6.4\times55\%=6.4\times(45\%+55\%)=6.4\times1=6.4$.
\dapso{$2.5$; $4.7$; $6.4$.}}
\gap
\begin{probbox}{13}{\relax}
Find $x$.
\begin{cols2}
\pa{a} $20\%\times x=1.35+0.25$ & \pa{b} $x:40\%-2.5=7.5:1.5$
\end{cols2}
\end{probbox}
\probsol{
\pa{a} $0.2\times x=1.6$, \ so $x=1.6:0.2=8$.\par
\pa{b} $x:40\%-2.5=5$, \ so $x:40\%=7.5$ \ and \ $x=7.5\times0.4=3$.
\dapso{\pa{a} $x=8$;\ \ \pa{b} $x=3$.}}
\gap
\begin{probbox}{14}{\relax}
Find $x$.
\begin{cols2}
\pa{a} $150\%\times x=4.2-1.8$ & \pa{b} $x:25\%+0.75=0.75\times21$
\end{cols2}
\end{probbox}
\probsol{
\pa{a} $1.5\times x=2.4$, \ so $x=2.4:1.5=1.6$.\par
\pa{b} $x:25\%+0.75=15.75$, \ so $x:25\%=15$ \ and \ $x=15\times0.25=3.75$.
\dapso{\pa{a} $x=1.6$;\ \ \pa{b} $x=3.75$.}}
\gap
\begin{probbox}{15}{\relax}
Find $x$.
\begin{cols2}
\pa{a} $x+35\%\times x=5.4$ & \pa{b} $x-12\%\times x=22$
\end{cols2}
\end{probbox}
\probsol{
\pa{a} $x=100\%\times x$, so $x+35\%\times x=135\%\times x=1.35\times x$.\par
$1.35\times x=5.4$, \ so $x=5.4:1.35=4$.\par
\pa{b} $x-12\%\times x=88\%\times x=0.88\times x=22$, \ so $x=22:0.88=25$.
\dapso{\pa{a} $x=4$;\ \ \pa{b} $x=25$.}}
\gap
\begin{probbox}{16}{\relax}
A tap fills an empty tank. In the first hour it fills $30\%$ of the tank; in the second hour it fills $12\%$ of the tank more than in the first hour. After these two hours, what percentage of the tank is still empty?
\end{probbox}
\probsol{
Second hour: $30\%+12\%=42\%$ of the tank.\par
Both hours: $30\%+42\%=72\%$ of the tank.\par
Still empty: $100\%-72\%=28\%$ of the tank.
\dapso{$28\%$ of the tank.}}
\gap
\begin{probbox}{17}{\relax}
A school library has three kinds of books: storybooks, science books and textbooks. Storybooks make up $45\%$ of all the books; the number of science books is $\dfrac23$ of the number of storybooks. What percentage of the books are textbooks?
\end{probbox}
\probsol{
Science books: $45\%\times\dfrac23=30\%$ of all the books.\par
Textbooks: $100\%-45\%-30\%=25\%$ of all the books.
\dapso{$25\%$.}}
\gap
\begin{probbox}{18}{\relax}
This year Khoa's age is $25\%$ of his father's age, and his grandfather's age is $175\%$ of his father's age. Khoa and his father are $50$ years old altogether. How old is each of the three people?
\end{probbox}
\probsol{
$25\%=\dfrac14$: Khoa's age is $1$ part, his father's age is $4$ parts.\par
Total number of parts: $1+4=5$; \ one part: $50:5=10$ (years).\par
Khoa: $10$ years; \ father: $10\times4=40$ years.\par
Grandfather: $40\times175\%=40\times1.75=70$ (years).
\dapso{Khoa $10$, father $40$, grandfather $70$ years old.}}
\gap
\begin{probbox}{19}{\relax}
A rectangle has a perimeter of $32$ dm. Its width is $60\%$ of its length. Find the area of the rectangle.
\end{probbox}
\probsol{
Length $+$ width: $32:2=16$ (dm).\par
$60\%=\dfrac35$: the width is $3$ parts, the length $5$ parts.\par
One part: $16:(3+5)=2$ (dm). \ Length: $2\times5=10$ dm; \ width: $2\times3=6$ dm.\par
Area: $10\times6=60$ (dm$^2$).
\dapso{$60$ dm$^2$.}}
\gap
\begin{probbox}{20}{\relax}
A farm has $150$ rabbits and chickens in all. After $10$ rabbits and $20$ chickens are sold, the number of rabbits left is $50\%$ of the number of chickens left. How many rabbits and how many chickens were there at first?
\end{probbox}
\probsol{
Animals left: $150-10-20=120$.\par
$50\%=\dfrac12$: rabbits left are $1$ part, chickens left are $2$ parts.\par
One part: $120:3=40$. \ Rabbits left: $40$; \ chickens left: $80$.\par
At first: rabbits $40+10=50$; \ chickens $80+20=100$.
\dapso{$50$ rabbits and $100$ chickens.}}
\gap
\begin{probbox}{21}{\relax}
Two shelves hold $72$ books altogether. If $20\%$ of the books on the first shelf are moved to the second shelf, the two shelves will hold the same number of books. How many books are on each shelf?
\end{probbox}
\probsol{
After the move each shelf holds $72:2=36$ books.\par
The first shelf keeps $100\%-20\%=80\%$ of its books, so $80\%$ of its books is $36$.\par
First shelf: $36:80\%=36:0.8=45$ (books).\par
Second shelf: $72-45=27$ (books).\par
\emph{Check:} $20\%$ of $45$ is $9$; \ $45-9=36$ and $27+9=36$.
\dapso{$45$ books and $27$ books.}}
\gap
\begin{probbox}{22}{\relax}
Two warehouses A and B store rice. If $2$ tonnes are moved from A to B, the two warehouses hold equal amounts. If instead $1$ tonne is moved from B to A, the rice in A is $300\%$ of the rice in B. How much rice did each warehouse hold at first?
\end{probbox}
\probsol{
Moving $2$ tonnes from A to B makes them equal, so A has $2\times2=4$ tonnes more than B.\par
If $1$ tonne goes from B to A, the difference becomes $4+1+1=6$ (tonnes).\par
Then A is $300\%=3$ times B: B is $1$ part, A is $3$ parts, the difference is $2$ parts.\par
One part: $6:2=3$ (tonnes). \ B then has $3$ tonnes, so at first B had $3+1=4$ tonnes.\par
At first A had $4+4=8$ tonnes.
\dapso{A $8$ tonnes, B $4$ tonnes.}}
\gap
\begin{probbox}{23}{\relax}
At the start of the school year, the number of boys in class 5A was $80\%$ of the number of girls. In the middle of the year, $2$ more boys joined the class, and then the number of boys was $90\%$ of the number of girls. How many pupils did class 5A have at the start of the year?
\end{probbox}
\probsol{
The number of girls did not change.\par
The $2$ new boys make up $90\%-80\%=10\%$ of the number of girls.\par
Girls: $2:10\%=2:0.1=20$.\par
Boys at the start: $20\times80\%=16$.\par
Pupils at the start: $20+16=36$.
\dapso{$36$ pupils.}}
\gap
\begin{probbox}{24}{\relax}
A shop had a stock of fabric. On Monday it sold $18.5\%$ of the stock, on Tuesday $26.5\%$, and on Wednesday as much as on Monday and Tuesday together. What percentage of the stock was left?
\end{probbox}
\probsol{
Monday and Tuesday: $18.5\%+26.5\%=45\%$ of the stock.\par
Wednesday: also $45\%$.\par
Left: $100\%-45\%-45\%=10\%$ of the stock.
\dapso{$10\%$.}}
\gap
\begin{probbox}{25}{\relax}
Find $x$.
\begin{cols2}
\pa{a} $75\%\times x+x\times\dfrac14=3.6$ & \pa{b} $x:20\%-x=12$
\end{cols2}
\end{probbox}
\probsol{
\pa{a} $\dfrac14=25\%$, so $75\%\times x+25\%\times x=100\%\times x=x$. \ Hence $x=3.6$.\par
\pa{b} $x:20\%=x:0.2=x\times5$, so $x\times5-x=x\times4=12$ \ and \ $x=3$.
\dapso{\pa{a} $x=3.6$;\ \ \pa{b} $x=3$.}}
\gap
\begin{probbox}{26}{\relax}
Calculate: $\ (45\%+0.55)\times3.7-75\%\times4$.
\end{probbox}
\probsol{
$45\%+0.55=0.45+0.55=1$, \ and $75\%\times4=0.75\times4=3$.\par
So the expression equals $1\times3.7-3=0.7$.
\dapso{$0.7$.}}
\gap
\begin{probbox}{27}{\relax}
In a school, $60\%$ of the pupils are girls, and $25\%$ of the girls sing in the school choir. What percentage of all the pupils are girls who sing in the choir?
\end{probbox}
\probsol{
The choir girls are $25\%$ of $60\%$ of the pupils:\par
$60\%\times25\%=0.6\times0.25=0.15=15\%$.
\dapso{$15\%$ of the pupils.}}
\gap
\begin{probbox}{28}{\relax}
The sum of two numbers is $2.7$. The first number is $125\%$ of the second. Find the two numbers.
\end{probbox}
\probsol{
$125\%=\dfrac{125}{100}=\dfrac54$: the first number is $5$ parts, the second $4$ parts.\par
One part: $2.7:(5+4)=0.3$.\par
First number: $0.3\times5=1.5$; \ second number: $0.3\times4=1.2$.
\dapso{$1.5$ and $1.2$.}}
\gap
\begin{probbox}{29}{\relax}
The difference between two numbers is $3.6$. The smaller number is $40\%$ of the larger number. Find the two numbers.
\end{probbox}
\probsol{
$40\%=\dfrac25$: the smaller number is $2$ parts, the larger $5$ parts.\par
The difference is $5-2=3$ parts, so one part is $3.6:3=1.2$.\par
Smaller: $1.2\times2=2.4$; \ larger: $1.2\times5=6$.
\dapso{$2.4$ and $6$.}}
\gap
\begin{probbox}{30}{\relax}
Mum cuts $35\%$ off a ribbon, then cuts $40\%$ of what is left. \ \pa{a} What percentage of the ribbon is left? \ \pa{b} The piece left is $7.8$ m long. How long was the ribbon?
\end{probbox}
\probsol{
\pa{a} After the first cut: $100\%-35\%=65\%$ of the ribbon is left.\par
The second cut leaves $100\%-40\%=60\%$ of that: $\ 65\%\times60\%=0.65\times0.6=0.39=39\%$.\par
\pa{b} $39\%$ of the ribbon is $7.8$ m, so the ribbon was $7.8:0.39=20$ (m).
\dapso{\pa{a} $39\%$;\ \ \pa{b} $20$ m.}}
\dividerbar{Challenge -- Extension}
\begin{probbox}{31}{\stars{1}}
At a swimming pool, the number of children sitting on the side is $25\%$ of the number of children in the water. Then $4$ children from the side jump into the water, and now the number of children on the side is $12.5\%$ of the number in the water. How many children are there in all?
\end{probbox}
\probsol{
$25\%=\dfrac14$: at first the children on the side are $\dfrac1{1+4}=\dfrac15$ of all the children.\par
$12.5\%=\dfrac18$: afterwards the children on the side are $\dfrac1{1+8}=\dfrac19$ of all the children.\par
The total does not change. The $4$ children who jumped in make up $\dfrac15-\dfrac19=\dfrac4{45}$ of the total.\par
Total: $4:\dfrac4{45}=45$ (children).\par
\emph{Check:} at first $9$ on the side, $36$ in the water ($9$ is $25\%$ of $36$); afterwards $5$ and $40$ ($5$ is $12.5\%$ of $40$).
\dapso{$45$ children.}}
\gap
\begin{probbox}{32}{\stars{1}}
An's stamps are $75\%$ of Bình's stamps. If each of them gets $10$ more stamps, An's stamps will be $80\%$ of Bình's. How many stamps does each boy have now?
\end{probbox}
\probsol{
$75\%=\dfrac34$: An has $3$ parts, Bình $4$ parts, so Bình has $1$ part more than An.\par
After each gets $10$ stamps, the difference is still $1$ part.\par
$80\%=\dfrac45$: then An has $4$ new parts and Bình $5$ new parts, a difference of $1$ new part.\par
So $1$ old part $=$ $1$ new part, and Bình's stamps go from $4$ parts to $5$ parts: the $10$ stamps are $1$ part.\par
Now: An $3\times10=30$ stamps, Bình $4\times10=40$ stamps.\par
\emph{Check:} $40:50=0.8=80\%$.
\dapso{An $30$ stamps, Bình $40$ stamps.}}
\gap
\begin{probbox}{33}{\stars{1}}
The price of a toy was raised by $20\%$. Later the new price was lowered by $20\%$. Compared with the original price, did the final price go up or down, and by what percentage?
\end{probbox}
\probsol{
After the rise the price is $120\%$ of the original.\par
After the cut it is $80\%$ of that: $\ 120\%\times80\%=1.2\times0.8=0.96=96\%$ of the original.\par
So the price went down by $100\%-96\%=4\%$.
\dapso{down by $4\%$.}}
\gap
\begin{probbox}{34}{\stars{1}}
The length of a rectangle is increased by $25\%$ and its width is decreased by $20\%$. How does the area change?
\end{probbox}
\probsol{
New length: $125\%$ of the old length; new width: $80\%$ of the old width.\par
New area: $125\%\times80\%=1.25\times0.8=1=100\%$ of the old area.
\dapso{the area does not change.}}
\gap
\begin{probbox}{35}{\stars{1}}
Three numbers add up to $108$. The first number is $50\%$ of the second, and the third is $150\%$ of the second. Find the three numbers.
\end{probbox}
\probsol{
$50\%=\dfrac12$ and $150\%=\dfrac32$. \ Take the second number as $2$ parts: the first is $1$ part, the third is $3$ parts.\par
One part: $108:(1+2+3)=18$.\par
The numbers: $18$, \ $18\times2=36$, \ $18\times3=54$.
\dapso{$18$; $36$; $54$.}}
\gap
\begin{probbox}{36}{\stars{1}}
Find $x$: $\ (x-15\%\times x):1.7=30\%$.
\end{probbox}
\probsol{
$x-15\%\times x=85\%\times x=0.85\times x$.\par
$0.85\times x=30\%\times1.7=0.3\times1.7=0.51$.\par
$x=0.51:0.85=0.6$.
\dapso{$x=0.6$.}}
\gap
\begin{probbox}{37}{\stars{1}}
Calculate: $\ A=1\%+2\%+3\%+\dots+98\%+99\%$. \ Write the answer as a decimal.
\end{probbox}
\probsol{
$A=(1+2+3+\dots+99)\%$.\par
$1+2+\dots+99=(1+99)\times99:2=4\,950$.\par
$A=4\,950\%=49.5$.
\dapso{$49.5$.}}
\gap
\begin{probbox}{38}{\stars{1}}
Calculate in a clever way: $\ 12.5\%\times0.25\times7.9\times8\times400\%$.
\end{probbox}
\probsol{
Group the factors: $(12.5\%\times8)\times(0.25\times400\%)\times7.9$.\par
$12.5\%\times8=0.125\times8=1$; \quad $0.25\times400\%=0.25\times4=1$.\par
Result: $1\times1\times7.9=7.9$.
\dapso{$7.9$.}}
\gap
\begin{probbox}{39}{\stars{1}}
A box contains red and blue balls; the red balls are $60\%$ of all the balls. After $20$ blue balls are put in, the red balls are $40\%$ of all the balls. How many red balls are in the box?
\end{probbox}
\probsol{
The number of red balls does not change.\par
At first: red $60\%$, blue $40\%$ of the balls, so blue $=\dfrac{40}{60}=\dfrac23$ of red.\par
Afterwards: red $40\%$, blue $60\%$, so blue $=\dfrac{60}{40}=\dfrac32$ of red.\par
The $20$ blue balls added are $\dfrac32-\dfrac23=\dfrac56$ of the red balls.\par
Red balls: $20:\dfrac56=24$.\par
\emph{Check:} at first $24$ red, $16$ blue ($24:40=60\%$); after: $24$ red, $36$ blue ($24:60=40\%$).
\dapso{$24$ red balls.}}
\gap
\begin{probbox}{40}{\stars{1}}
Three taps fill a tank together. Tap A pours in $40\%$ of the tank, tap B pours in $75\%$ of what tap A pours, and tap C pours in the rest. Tap C pours $54$ litres less than tap A. How many litres does the tank hold?
\end{probbox}
\probsol{
Tap B: $40\%\times75\%=0.4\times0.75=0.3=30\%$ of the tank.\par
Tap C: $100\%-40\%-30\%=30\%$ of the tank.\par
Tap A pours $40\%-30\%=10\%$ of the tank more than tap C, and this is $54$ litres.\par
The tank: $54:10\%=54:0.1=540$ (litres).
\dapso{$540$ litres.}}
\gap
\begin{probbox}{41}{\stars{1}}
A shop sold $40\%$ of its rice on the first day and $50\%$ of the remaining rice on the second day; $48$ kg were left. How many kilograms of rice did the shop have?
\end{probbox}
\probsol{
\emph{Work backwards.} After the second day, $100\%-50\%=50\%$ of the remainder is left, and this is $48$ kg.\par
Remainder after day one: $48:50\%=96$ (kg).\par
This is $100\%-40\%=60\%$ of all the rice, so the shop had $96:60\%=96:0.6=160$ (kg).
\dapso{$160$ kg.}}
\gap
\begin{probbox}{42}{\stars{1}}
Find $x$: $\ x\times125\%+x:4+x\times0.5=3.4$.
\end{probbox}
\probsol{
$x:4=x\times0.25$, so the left side is $x\times(1.25+0.25+0.5)=x\times2$.\par
$x\times2=3.4$, \ so $x=1.7$.
\dapso{$x=1.7$.}}
\gap
\begin{probbox}{43}{\stars{1}}
This year a daughter's age is $25\%$ of her mother's age. In $6$ years' time, the daughter's age will be $\dfrac13$ of her mother's age. How old is each of them now?
\end{probbox}
\probsol{
Now: daughter $1$ part, mother $4$ parts; the difference is $3$ parts.\par
In $6$ years: daughter $1$ new part, mother $3$ new parts; the difference is $2$ new parts.\par
The age difference never changes, so $3$ parts $=2$ new parts, i.e.\par $1$ new part $=1.5$ parts.\par
The daughter grows from $1$ part to $1.5$ parts: $0.5$ part $=6$ years, so $1$ part $=12$ years.\par
Daughter: $12$ years; \ mother: $12\times4=48$ years.\par
\emph{Check:} in $6$ years, $18$ and $54$, and $18=\dfrac13$ of $54$.
\dapso{daughter $12$, mother $48$ years old.}}
\gap
\begin{probbox}{44}{\stars{1}}
A jug holds $200$ g of salt water in which salt makes up $5\%$. How many grams of pure water must be added so that salt makes up $4\%$?
\end{probbox}
\probsol{
Salt in the jug: $200\times5\%=10$ (g). Adding water does not change this.\par
After adding, $10$ g is $4\%$ of the salt water: $\ 10:4\%=10:0.04=250$ (g).\par
Water to add: $250-200=50$ (g).
\dapso{$50$ g of water.}}
\gap
\begin{probbox}{45}{\stars{1}}
Number A is $25\%$ more than number B. By what percentage is B smaller than A?
\end{probbox}
\probsol{
A is $125\%$ of B, i.e.\par $A=\dfrac54$ of B: B is $4$ parts, A is $5$ parts.\par
B is smaller by $1$ part, and $1:5=0.2=20\%$ of A.
\dapso{$20\%$.}}
\gap
\begin{probbox}{46}{\stars{1}}
In a class, $45\%$ of the pupils are boys, and there are $6$ more girls than boys. How many pupils are in the class?
\end{probbox}
\probsol{
Girls: $100\%-45\%=55\%$ of the class.\par
The difference: $55\%-45\%=10\%$ of the class, which is $6$ pupils.\par
The class: $6:10\%=60$ (pupils).
\dapso{$60$ pupils.}}
\gap
\begin{probbox}{47}{\stars{2}}
Fresh mushrooms contain $92\%$ water. After drying, the mushrooms contain $20\%$ water. How many kilograms of dried mushrooms do we get from $100$ kg of fresh mushrooms?
\end{probbox}
\probsol{
In $100$ kg of fresh mushrooms the dry matter is $100\%-92\%=8\%$: $\ 100\times8\%=8$ (kg).\par
Drying only removes water, so the dried mushrooms still contain $8$ kg of dry matter.\par
In dried mushrooms the dry matter is $100\%-20\%=80\%$.\par
Dried mushrooms: $8:80\%=8:0.8=10$ (kg).
\dapso{$10$ kg.}}
\gap
\begin{probbox}{48}{\stars{1}}
The sum of two numbers is $168$; $40\%$ of the first number equals $30\%$ of the second number. Find the two numbers.
\end{probbox}
\probsol{
$\dfrac4{10}$ of the first number $=\dfrac3{10}$ of the second, so $4$ times the first $=3$ times the second.\par
Hence the first number is $3$ parts and the second $4$ parts.\par
One part: $168:7=24$. \ The numbers: $24\times3=72$ and $24\times4=96$.\par
\emph{Check:} $40\%\times72=28.8=30\%\times96$.
\dapso{$72$ and $96$.}}
\gap
\begin{probbox}{49}{\stars{2}}
A number is decreased by $25\%$. By what percentage must the new number be increased to get back to the original number?
\end{probbox}
\probsol{
Let the original number be $4$ parts. Decreasing it by $25\%=\dfrac14$ leaves $3$ parts.\par
To get back we must add $1$ part, which is $\dfrac13$ of the new number.\par
$\dfrac13=1:3=33\dfrac13\%$.
\dapso{$33\dfrac13\%$.}}
\gap
\begin{probbox}{50}{\stars{1}}
An, Bình and Chi share the cost of a present. An pays $30\%$ of the cost; Bình pays $40\%$ of the remaining amount; Chi pays the rest. Chi pays $14\,000$ dong more than Bình. How much does each person pay?
\end{probbox}
\probsol{
After An: $100\%-30\%=70\%$ of the cost remains.\par
Bình: $70\%\times40\%=28\%$ of the cost; \ Chi: $70\%-28\%=42\%$ of the cost.\par
Chi pays $42\%-28\%=14\%$ of the cost more than Bình, and this is $14\,000$ dong.\par
Cost: $14\,000:14\%=14\,000:0.14=100\,000$ (dong).\par
An: $30\,000$ dong; \ Bình: $28\,000$ dong; \ Chi: $42\,000$ dong.
\dapso{An $30\,000$ dong, Bình $28\,000$ dong, Chi $42\,000$ dong.}}
